% Pencil scenes are illustrative teaching aids. TeX owns the role labels.

\tikzset{
  sung role tag/.style={
    draw=black,fill=white,line width=0.5pt,rounded corners=1.3pt,
    inner xsep=2.5pt,inner ysep=1.2pt,
    font=\fontsize{7.8}{9}\selectfont\bfseries}}

\newcommand{\SungRoleTag}[4]{%
  \node[sung role tag,anchor=south]
    at ($(#1.south west)!#2!(#1.south east)+(0,#3)$) {#4};}

\newcommand{\SungAboveTag}[3]{%
  \node[sung role tag,anchor=south] at
    ($(#1.north west)!#2!(#1.north east)+(0,1mm)$) {#3};}
\newcommand{\SungBelowTag}[3]{%
  \node[sung role tag,anchor=north] at
    ($(#1.south west)!#2!(#1.south east)+(0,-1mm)$) {#3};}

\newcommand{\CommunionPatenSafetyDiagram}{%
  \begin{center}
  \begin{tikzpicture}
    \node[inner sep=0pt] (scene)
      {\includegraphics[width=.40\linewidth]{%
        liturgy/roman-rite/1962/reference/altar-server-guides/shared/artwork/pencil/ASG-ART-026-communion-paten-overhead-v1-alpha.png}};
    \SungRoleTag{scene}{.25}{27mm}{PRIEST'S HAND}
    \SungRoleTag{scene}{.58}{27mm}{COMMUNICANT}
    \node[sung role tag,anchor=south] at
      ($(scene.south west)!.69!(scene.south east)+(0,1mm)$)
      {PLATE LEVEL UNDER CHIN};
    \node[sung role tag,anchor=north west,align=left] at
      ($(scene.south west)!.02!(scene.south east)+(0,1mm)$)
      {SIDE PROFILE\\FOLLOW THE HOST WITHOUT TOUCHING\\THE COMMUNICANT};
  \end{tikzpicture}
  \end{center}}

\newcommand{\MCGospelFormationDiagram}{%
  \begin{center}\begin{tikzpicture}
    \node[inner sep=0pt] (scene) {\includegraphics[width=.90\linewidth]{%
      liturgy/roman-rite/1962/reference/altar-server-guides/shared/artwork/pencil/ASG-ART-009-mc-gospel-formation-2026.png}};
    \SungAboveTag{scene}{.164}{P}
    \SungBelowTag{scene}{.105}{A1}
    \SungBelowTag{scene}{.261}{A2}
    \SungBelowTag{scene}{.305}{MC}
  \end{tikzpicture}

  \parbox{.90\linewidth}{\small\textbf{Stationary Gospel service, without incense.}
  The Missal remains on its altar stand at the Gospel side. P chants at the
  book; A1 and A2 remain at their assigned candle stations; MC attends beside
  the book. This model selects six altar candles and local predella places. On the incense route,
  TH also waits outward with the closed thurible. Mark the actual distances
  and route before rehearsal.}
  \end{center}}

\newcommand{\MCEntranceHandoffsDiagram}{%
  \begin{center}
  \begin{tikzpicture}
    \node[inner sep=0pt] (scene)
      {\includegraphics[width=.95\linewidth]{%
        liturgy/roman-rite/1962/reference/altar-server-guides/shared/artwork/pencil/ASG-ART-022-mc-entrance-handoffs-v2-alpha.png}};
    \SungRoleTag{scene}{.074}{58mm}{ADULT}
    \SungRoleTag{scene}{.223}{58mm}{A1}
    \SungRoleTag{scene}{.337}{58mm}{A2}
    \SungRoleTag{scene}{.414}{66mm}{P}
    \SungRoleTag{scene}{.462}{58mm}{MC}
    \SungRoleTag{scene}{.590}{58mm}{A1}
    \SungRoleTag{scene}{.703}{58mm}{P}
    \SungRoleTag{scene}{.838}{58mm}{MC}
    \SungRoleTag{scene}{.952}{58mm}{A2}
  \end{tikzpicture}

  \parbox{.95\linewidth}{\small\textbf{Two controlled handoffs.}
  \textit{Before the entrance:} the adult lights and gives one candle to each
  acolyte; MC remains free-handed and the priest retains his biretta.
  \textit{After the objects are placed at their named local stations:} MC
  receives the biretta two-handed while A1 and A2 wait empty-handed and keep
  the centre clear. The drawing does not prescribe the procession order,
  arrival reverence, distances, or church layout; mark those locally before
  rehearsal.}
  \end{center}}

\newcommand{\MCElevationTableauDiagram}{%
  \begin{center}
  \begin{tikzpicture}
    \node[inner sep=0pt] (scene)
      {\includegraphics[width=.82\linewidth]{%
        liturgy/roman-rite/1962/reference/altar-server-guides/shared/artwork/pencil/ASG-ART-024-mc-elevation-tableau-v2-alpha.png}};
    \SungRoleTag{scene}{.150}{34mm}{T1}
    \SungRoleTag{scene}{.337}{34mm}{T2}
    \SungRoleTag{scene}{.525}{34mm}{T3}
    \SungRoleTag{scene}{.703}{34mm}{T4}
    \SungRoleTag{scene}{.891}{29mm}{A1}
  \end{tikzpicture}

  \parbox{.82\linewidth}{\small\textbf{Stable elevation service
  (MC-037--041).} T1--T4 kneel at the four locally marked places and hold
  their torches upright with both hands. A1 kneels separately, gives only the
  controlled bell signals at MC-039, and replaces the bell on its pad after
  the chalice-elevation group. The drawing fixes neither spacing nor
  architecture and does not show the optional incense branch.}
  \end{center}}

\newcommand{\SolemnGospelStationDiagram}{%
  \begin{center}\begin{tikzpicture}
    \node[inner sep=0pt] (scene) {\includegraphics[width=.62\linewidth]{%
      liturgy/roman-rite/1962/reference/altar-server-guides/shared/artwork/pencil/ASG-ART-010-so-gospel-station-2026.png}};
    \SungAboveTag{scene}{.22}{A1}
    \SungAboveTag{scene}{.40}{SD}
    \SungAboveTag{scene}{.63}{A2}
    \SungBelowTag{scene}{.21}{TH}
    \SungBelowTag{scene}{.56}{D}
    \SungBelowTag{scene}{.79}{MC}
  \end{tikzpicture}

  \parbox{.90\linewidth}{\small \textbf{Stationary Gospel group.} SD faces D and supports the open book between A1 and A2. D chants with joined hands; MC attends at D\textquotesingle s right and TH waits with the closed thurible at his left. This close view identifies the roles and their relation to the book; mark distances and the route locally.}
  \end{center}}

\newcommand{\SolemnGospelDepartureDiagram}{%
  \begin{center}\begin{tikzpicture}
    \node[inner sep=0pt] (scene) {\includegraphics[width=.90\linewidth]{%
      liturgy/roman-rite/1962/reference/altar-server-guides/shared/artwork/pencil/ASG-ART-025-so-gospel-departure-2026.png}};
    \SungBelowTag{scene}{.12}{MC}
    \SungBelowTag{scene}{.27}{TH}
    \SungBelowTag{scene}{.41}{A1}
    \SungBelowTag{scene}{.55}{A2}
    \SungBelowTag{scene}{.70}{SD}
    \SungBelowTag{scene}{.87}{D}
  \end{tikzpicture}

  \parbox{.90\linewidth}{\small \textbf{Single-file Gospel departure (SO-020--021).} MC leads; TH follows with the closed thurible; A1 and A2 follow in successive places with lighted candles; SD follows empty-handed, then D carries the closed Gospel book with both hands. Bb has returned the boat and remains at base. This shows the model order; mark the route, distances, and common-reverence place locally.}
  \end{center}}

\newcommand{\SolemnEntranceHandoffsDiagram}{%
  \begin{center}\begin{tikzpicture}
    \node[inner sep=0pt] (scene) {\includegraphics[width=.95\linewidth]{%
      liturgy/roman-rite/1962/reference/altar-server-guides/shared/artwork/pencil/ASG-ART-023-so-entrance-handoffs-2026.png}};
    \SungAboveTag{scene}{.10}{ADULT}
    \SungAboveTag{scene}{.21}{A1}
    \SungAboveTag{scene}{.29}{A2}
    \SungAboveTag{scene}{.49}{C}
    \SungAboveTag{scene}{.63}{D}
    \SungAboveTag{scene}{.75}{SD}
    \SungAboveTag{scene}{.92}{MC}
  \end{tikzpicture}

  \parbox{.90\linewidth}{\small \textbf{Two separate handoff moments.} Before the entrance, the adult gives a lighted candle to each acolyte. After the sacred ministers uncover at the sanctuary, MC receives their birettas one at a time with both hands and places each at its marked sedilia place. Here SD hands his to MC while C and D retain theirs. This is a handoff study; rehearse procession order, arrival reverence, and distances at the actual church.}
  \end{center}}

\newcommand{\MCFirstAblutionDiagram}{%
  \begin{center}
  \begin{tikzpicture}
    \node[inner sep=0pt] (scene)
      {\includegraphics[width=.66\linewidth]{%
        liturgy/roman-rite/1962/reference/altar-server-guides/shared/artwork/pencil/ASG-ART-017-lm-mc-first-ablution-wine-v2-alpha.png}};
    \SungRoleTag{scene}{.18}{1mm}{A1: WINE}
    \SungRoleTag{scene}{.50}{40mm}{P}
    \SungRoleTag{scene}{.82}{1mm}{A2: WATER WAITS}
  \end{tikzpicture}

  \parbox{.82\linewidth}{\small\textbf{First ablution (MC-048; MODEL).}
  A1 pours wine into the priest-held chalice. A2 waits with water at the
  locally marked handoff point.}
  \end{center}}

\newcommand{\MCFingerAblutionDiagram}{%
  \begin{center}\begin{tikzpicture}
    \node[inner sep=0pt] (scene) {\includegraphics[width=.90\linewidth]{%
      liturgy/roman-rite/1962/reference/altar-server-guides/shared/artwork/pencil/ASG-ART-018-mc-finger-ablutions-2026.png}};
    \SungAboveTag{scene}{.23}{P}
    \SungAboveTag{scene}{.75}{P}
    \SungBelowTag{scene}{.08}{A1: WINE}
    \SungBelowTag{scene}{.40}{A2: WAITS}
    \SungBelowTag{scene}{.59}{A1: DONE}
    \SungBelowTag{scene}{.92}{A2: WATER}
  \end{tikzpicture}

  \parbox{.90\linewidth}{\small \textbf{Finger ablution (MC-049; MODEL), two successive moments.} Left: A1 pours wine over the priest\textquotesingle s thumbs and index fingers into the chalice while A2 waits. Right: A1 has stopped; A2 then pours water over the same fingertips. Mark standing places and pouring signals locally.}
  \end{center}}

\newcommand{\MCAblutionsDiagram}{%
  \MCFirstAblutionDiagram
  \MCFingerAblutionDiagram}

\newcommand{\SolemnAblutionsDiagram}{%
  \par\noindent\begin{minipage}{\linewidth}
  \begin{center}\begin{tikzpicture}
    \node[inner sep=0pt] (scene)
      {\includegraphics[width=.96\linewidth]{%
        liturgy/roman-rite/1962/reference/altar-server-guides/shared/artwork/pencil/ASG-ART-030-so-subdeacon-ablutions-2026.png}};
    \SungAboveTag{scene}{.167}{FIRST WINE}
    \SungAboveTag{scene}{.50}{FINGER ABLUTION: WINE}
    \SungAboveTag{scene}{.833}{THEN WATER}
    \SungBelowTag{scene}{.05}{A1}
    \SungBelowTag{scene}{.17}{C}
    \SungBelowTag{scene}{.28}{SD}
    \SungBelowTag{scene}{.38}{A1}
    \SungBelowTag{scene}{.50}{C}
    \SungBelowTag{scene}{.61}{SD}
    \SungBelowTag{scene}{.71}{A1}
    \SungBelowTag{scene}{.84}{C}
    \SungBelowTag{scene}{.94}{SD}
  \end{tikzpicture}

  \parbox{.96\linewidth}{\small\textbf{Three successive Solemn-Mass moments
  (SO-057--058; MODEL), hand-action close-ups.} A1 supplies and receives both
  cruets at the LOCAL handoff point. The subdeacon pours the first wine directly
  into the priest-held chalice. He then pours wine, followed by water, over the
  celebrant's thumbs and index fingers into the chalice.}
  \end{center}\end{minipage}\par}
